% Commutative diagram from the long-exact-sequence chapter.
% Source: book/tex/homology/long-exact.tex (lines 27-45).
\documentclass[border=2mm]{standalone}
\usepackage{amssymb,amsmath}
\usepackage{tikz-cd}
\newcommand{\ZZ}{\mathbb{Z}}
\newcommand{\eps}{\varepsilon}
\newcommand{\id}{\mathrm{id}}
\begin{document}
\begin{tikzcd}[column sep=huge]
	& \vdots \ar[d, "\partial_A"] & \vdots \ar[d, "\partial_B"] & \vdots \ar[d, "\partial_C"] & \\
	0 \ar[r]
		& A_{n+1} \ar[hook, r, "f_{n+1}"] \ar[d, "\partial_A"]
		& B_{n+1} \ar[r, two heads, "g_{n+1}"] \ar[d, "\partial_B"]
		& C_{n+1} \ar[r] \ar[d, "\partial_C"]
		& 0 \\
	0 \ar[r]
		& A_n \ar[hook, r, "f_n"] \ar[d, "\partial_A"]
		& B_n \ar[r, two heads, "g_n"] \ar[d, "\partial_B"]
		& C_n \ar[r] \ar[d, "\partial_C"]
		& 0 \\
	0 \ar[r]
		& A_{n-1} \ar[hook, r, "f_{n-1}"] \ar[d, "\partial_A"]
		& B_{n-1} \ar[r, two heads, "g_{n-1}"] \ar[d, "\partial_B"]
		& C_{n-1} \ar[r] \ar[d, "\partial_C"]
		& 0 \\
	& \vdots & \vdots & \vdots
\end{tikzcd}
\end{document}
